\documentclass[10pt,a4paper]{amsart}
\usepackage[margin=19mm]{geometry}
\usepackage{amsmath,amssymb,mathtools}
\usepackage[scr=rsfs,cal=euler]{mathalfa}
\usepackage{xfrac}
\usepackage[unicode,hidelinks,bookmarks=false]{hyperref}
\newcommand{\ie}{i.e.\ }
\newcommand{\cf}{cf.\ }
\newcommand{\resp}{respectively\ }
\newcommand{\Subsection}{\subsection}
\providecommand{\nobreakdash}{\nobreak\hbox{-}}
\setlength{\emergencystretch}{2em}
\multlinegap0pt
\theoremstyle{plain}
\newtheorem{thm}{}
\title[Excluding a fat tree]{Asymptotic structure. III. Excluding a fat tree: a geodesic subdivision of the fat tree forces it as a fat minor}
\author{Tung Nguyen, Alex Scott, Paul Seymour}
\date{Source: arXiv:2509.09035v4 (CC BY 4.0)}
\newcommand\blackslug{\hbox{\hskip 1pt \vrule width 4pt height 8pt depth 1.5pt \hskip 1pt}}
\newcommand\bbox{\hfill \quad \blackslug \bigbreak}
\def\LL{,\ldots,}
\newcommand{\voro}{\operatorname{\Delta}}
\newcommand{\bd}{\operatorname{bd}}
\newcommand{\mac}{\mathcal}
\newcommand{\Proof}{\noindent{\bf Proof.}\ \ }
\ifdefined\LL\else
\def\LL{,\ldots,}
\fi
\ifdefined\Proof\else
\newcommand{\Proof}{\noindent{\bf Proof.}\ \ }
\fi
\ifdefined\bd\else
\newcommand{\bd}{\operatorname{bd}}
\fi
\ifdefined\mac\else
\newcommand{\mac}{\mathcal}
\fi
\ifdefined\voro\else
\newcommand{\voro}{\operatorname{\Delta}}
\fi
\begin{document}
\maketitle
\noindent\emph{Source and licence.} Asymptotic structure. III. Excluding a fat tree. Version arXiv:2509.09035v4 (CC BY 4.0), distributed under the Creative Commons Attribution 4.0 International licence, \url{http://creativecommons.org/licenses/by/4.0/}, as recorded in the arXiv licence metadata for this exact version and shown on the abstract page \url{https://arxiv.org/abs/2509.09035v4}. Full licence notice: licenses/arxiv-2509.09035-CC-BY-4.0.txt.\par\smallskip

\noindent\textit{Editorial scope.} Counted unit: the article's Theorem with source label \texttt{notmuchleft} (source0.tex lines 2579--2586), which bounds the quasi-bound of the Voronoi cell of a maximal adjoining-connected set of rebels with respect to a stable government, by an explicit pair depending on the level, the fatness parameter and the two canonicity parameters; delivered together with the article's complete original proof of it (source0.tex lines 2587--2700), which the article writes as a run-in proof block introduced by its own proof macro and closes with its proof-end mark. No mathematics is written, repaired or re-derived: statement and proof are the article's own text line for line, in the article's own notation (governments, rebels, adjoining-connectedness, Voronoi cells, quasi-bounds and the canonicity parameters), and no retained line is substituted, re-flowed or omitted, including the commented-out display carried verbatim inside the counted statement. Required context retained uncounted, in source order: the statements of the four results the counted proof invokes, namely \texttt{contactgraph} (lines 1809--1830), \texttt{smallcoarsepw} (lines 1898--1905), \texttt{palacebound} (lines 1986--1992) and \texttt{cabalsaresmall} (lines 2171--2175). The proofs of those context results are not delivered. Nothing else of the article is retained: its main theorems, its other lemmas and its figures are omitted. No citation command occurs in any retained line, so the article's bibliography is not delivered and the wrapper carries no bibliography entry, although the article ships one. Every cross-reference inside the retained spans resolves inside the delivered document. Editorial formatting only: an AMS \texttt{amsart} preamble that restates the environments and the macros the retained lines use, namely the article's declaration of its theorem environment with the article's own printed name (the article prints its theorems with the number only), its proof macro and proof-end mark, and the macros for the special letters, operators and decorations of the retained lines. The article is typeset in its own \texttt{article}-class format with its own fonts; no class or style file is shipped with this package and the wrapper is an amsart document. Numbering differs from the article, because the article numbers its theorems inside each section with a printed name consisting of the number only, whereas the wrapper keeps a continuous counter and retains five environments: the four retained context theorems print with the wrapper's own numbering in source order and the counted theorem as Theorem 5. One bundled source file, \texttt{sources/184869/source0.tex}, accompanies the delivered item as the exact source of every retained span. Every retained span is recorded in \texttt{delivery-evidence.json} with method \texttt{exact\_tex}.
	\begin{thm}\label{contactgraph}
		Let $\sigma$ be a standard, let
		$G$ be a graph, and
		let $\mathcal{T}$ be a $k$th-century realm in $G$, optimal for $\sigma$. 
		Each $\mac T$-community adjoins at most two forts, and 
		each fort $\mac T$-semiadjoins at most two other forts. Consequently, if 
		$C$ is an adjoin-connected subset of $\mac T^0$ that contains a fort, then  
		the forts in $C$ can be numbered as $X_i\;(i\in I)$, where $I$ is an integer interval, with the following properties:
		\begin{itemize}
			\item for each $i\in I$, if $i+1\in I$ then $X_i$ $C$-semiadjoins $X_{i+1}$, and no other pair of forts in $C$  $\mac T$-semiadjoin  each other 
			except possibly $X_a,X_b$, where $(a,b)$ is an end-set of $I$ and $b\ge a+2$;
			\item each $\mac T$-community  $S$  with $S\subseteq C$ adjoins at most two forts in $C$ (which therefore $C$-semiadjoin if there are two).
		\end{itemize}
		Moreover, if $C,D$ are disjoint adjoin-connected subsets of $\mac T^0$, and $P$ is a passage joining some $X\in C$ and some  $Y\in D$,
		then:
		\begin{itemize}
			\item  either $X$ is a $C$-peripheral fort, or $X$ belongs to a $C$-peripheral $C$-village, or $D$ contains no forts; and 
			\item either $Y$ 
			is a $D$-peripheral fort, or $Y$ belongs to a $D$-peripheral $D$-village,
			or $C$ contains no forts.
		\end{itemize}
	\end{thm}

	\begin{thm}\label{smallcoarsepw}
		Let $\sigma$ be a standard, let
		$G$ be a graph,
		let $\mathcal{T}$ be a $k$th-century realm in $G$, optimal for $\sigma$.
		Then $\voro_{\mac T} (\mac T^0)$
		has quasi-line-width at most
		$(10\alpha,\beta+d_0)$.
	\end{thm}

	\begin{thm}\label{palacebound}
		In the same notation, suppose that $G$ does not contain $H_\ell$ as a $c$-fat minor. Let $C$ be a $\mac T$-adjoin-connected set of rebels that contains at least
		one fort, and let $A\in \mac G$ be a palace.
		Let $U$ be the set of vertices in $\voro_{\mac G}(C)$ with a neighbour in  $\voro_{\mac G}(A)$.
		Then $U$ has quasi-size at most
		$$\left(4\cdot 3^{\ell-k-1}\alpha-2\alpha, \beta+4d_0+1\right).$$
	\end{thm}

	\begin{thm}\label{cabalsaresmall}
		With notation as above, suppose that $C$ is a cabal.
		Then $\bd(\voro_{\mac G}(C))$ has quasi-size at most
		$$(16(\ell-k-1)3^{\ell-k-1}\alpha, \beta+4d_0+1).$$
	\end{thm}

	\begin{thm}\label{notmuchleft}
		With notation as before, 
		let $\mac G$ be a stable government for $\mac T$, and
		let $C$ be a maximal $\mac G$-adjoin-connected set of rebels. Then
		$\voro_{\mac G}(C)$ has quasi-bound at most 
		$$\left(\ell\,3^{\ell+1}\alpha, \beta+4d_0+1\right).$$
		%$$\left(8(\ell-k-1) 3^{\ell-k-1}\alpha, \beta+4d_0+1\right).$$
	\end{thm}

	\Proof
	By \ref{smallcoarsepw}, $\voro_{\mac G}(C)\subseteq \voro_{\mac T}(C)$ has quasi-line-width at most $(10\alpha,\beta+d_0)$, so it remains to bound the quasi-size of $\bd(\voro_{\mac G}(C))$. If $k=\ell-1$, then there are no castles in $\mac T$ and no palaces in $\mac G$ 
	(because $G$ does not contain $H_{\ell}$ as a $c$-fat minor), and so $\bd(\voro_{\mac G}(C))=\emptyset$
	and the claim is true. So we may assume that $k\le \ell-2$.
	
	If $C$ is a $\mac T$-community, then, since $\bd(V(C))$ has quasi-size at most $(\alpha,\beta)$ and every vertex in $\bd(\voro_{\mac G}(C))$ has distance at most
	$d_0$ from $\bd(V(C))$, it follows that $\bd(\voro_{\mac G}(C))$ has quasi-size at most $(\alpha,\beta+d_0)$.
	So we may assume that $C$ is not a $\mac T$-community.
	
	We say that $C'\subseteq C$ is {\em dangerous} if $C'$ is $\mac G$-adjoin-connected, and 
	there exist $j>k$ and three members $A_1,A_2,A_3\in \mac G$, all of rank $j$, such that for $1\le i\le 3$, some member of $C'$ $\mac G$-adjoins $A_i$.
	If $C$ is dangerous, we want to choose a minimal subset $D'$ of $C$ that is still dangerous, and $\mac G$-village-closed, and $\mac G$-adjoin-connected, 
	\\
	\\
	(1) {\em $C$ is not dangerous.}
	\\
	\\
	Suppose that $C$ is dangerous, but let us digress for a moment and consider what we need. What we would really like is to obtain 
	a ``good'' dangerous subset $C'$, where ``good'' means 
	\begin{itemize}
		\item $C'$ is $\mac G$-village-closed and  $\mac G$-adjoin-connected, 
		\item only a bounded number of palaces in the government $\mac G$-adjoin members of $C'$;
		\item no rebel house or fort that is not in $C'$ $\mac G$-adjoins a non-$C'$-peripheral fort in $C$, or a member of a non-$C'$-peripheral $C'$-village; and
		\item the number of $C'$-peripheral $C'$-villages is bounded. (There are automatically at most two $C'$-peripheral forts.)
	\end{itemize}
	The third condition is needed for \ref{cabalsaresmall}, and the fourth condition is needed to prove \ref{palacebound}.
	
	$C$ itself satisfies the first and third conditions. By simply choosing $C'$ to be a minimal subset of $C$ that
	is $\mac G$-village-closed,  $\mac G$-adjoin-connected and dangerous, we can arrange the second bullet and fourth bullets, but that might wreck the third,
	so we need to be more cautious. Indeed, there are some cases where no good subset exists. For instance, if $C$ consists of a 
	single fort or a single $\mac G$-village, and $\mac G$-adjoins many palaces in $\mac G$, there is nothing we can do; or if $|C|=6$,
	and consists of three forts and three $\mac G$-villages making a six-cycle under $\mac G$-adjoinment, and $C$ is minimally dangerous,
	then again there is no good subset. This explains why cabals sometimes have ``leader'' terms that are not peripheral. 
	
	Let us continue the proof of (1).
	If some $\mac G$-village included in $C$ is dangerous, then it is a cabal (designating itself as the only leader $\mac G$-village), a contradiction. So
	no $\mac G$-village included in $C$ is dangerous.
	Thus we assume (for a contradiction) that $C$ is dangerous and not a $\mac T$-community, and no $\mac G$-village in $C$ is dangerous.
	Since $C$ is $\mac G$-adjoin-connected, it is $\mac T$-adjoin-connected, and so we may number the forts in
	$C$ as $X_i\;(i\in I)$ where $I$ is an integer interval, as in \ref{contactgraph}.
	Since $C$ is a maximal $\mac G$-adjoin-connected set of rebels, it follows that $C$ is $\mac G$-village-closed.
	
	For $i_1,i_2\in I$ with $i_1\le i_2$, let $C(i_1,i_2)$ be the union of $\{X_{i_1}\LL X_{i_2}\}$ and all $\mac G$-villages that $\mac G$-adjoin
	one of $X_{i_1}\LL X_{i_2}$.
	It follows that
	$C(i_1,i_2)$ is $\mac G$-village-closed.
	Since $C$ is dangerous, we may choose $i_1,i_2$ with $i_2-i_1$ minimal such that $C(i_1,i_2)$ is dangerous. (Possibly $i_1=i_2$.)
	
	Define $D$ as follows:
	\begin{itemize}
		\item if $i_2=i_1$, let $D=\{X_{i_1}\}$;
		\item if $i_2=i_1+1$, let $D$ be the union of $\{X_{i_1}, X_{i_2}\}$ and all $\mac G$-villages that $\mac G$-adjoin  both of $X_{i_1}, X_{i_2}$;
		\item if $i_2\ge i_1+2$, let 
		$D$ be the union of $\{X_{i_1}\LL X_{i_2}\}$ and all $\mac G$-villages that $\mac G$-adjoin one of $X_{i_1+1}\LL X_{i_2-1}$.
	\end{itemize}
	In each case, 
	$D$ is $\mac G$-village-closed, and $\mac G$-adjoin-connected.
	Each member of $C(i_1,i_2)\setminus D$ is a house, and belongs to a $\mac G$-village that $\mac G$-adjoins one or both of $X_{i_1},X_{i_2}$ and none of
	$X_{i_1+1}\LL X_{i_2-1}$.
	
	Since $C(i_1,i_2)$ is dangerous, there exist $j>k$ and three members $A_1,A_2,A_3\in \mac G$, all of rank $j$, such that for $1\le i\le 3$, some member of $C(i_1,i_2)$
	$\mac G$-adjoins  $A_i$.
	For $i = 1,2,3$, $A_i$ $\mac G$-adjoins either some member of $D$, or some $\mac G$-village
	included in $C(i_1,i_2)\setminus D$.
	So by adding to $D$ at most three of the $\mac G$-villages included in $C(i_1,i_2)\setminus D$, we can construct  a set $D'$ that is dangerous,
	$\mac G$-adjoin-connected, and $\mac G$-village-closed. Let us construct such a set $D'$ by adding to $D$ as few $\mac G$-villages as possible;  so in particular, if $D$ is dangerous then $D'=D$.
	We designate $X_{i_1}, X_{i_2}$ as leader forts of $D'$, and the (at most three) $\mac G$-villages with union $D'\setminus D$ as leader $\mac G$-villages of $D'$.
	We claim that $D'$ is a cabal (which will be a contradiction). To show this, we must check that
	\begin{itemize}
		\item $D'$ is $\mac G$-adjoin-connected, and $\mac G$-village-closed;
		\item at most two forts of $D'$ are designated as leader forts, and at most three $\mac G$-villages of $D'$ are designated as
		leader $\mac G$-villages;
		\item every $D'$-peripheral fort is a leader fort, and every $D'$-peripheral $D'$-village is a leader $\mac G$-village;
		\item if $X\in D'$ $\mac G$-adjoins some rebel not in $D'$, then either $X$ is a leader fort, or $X$ belongs to a leader $\mac G$-village;
		\item for some $j\in \{k+1\LL \ell-1\}$ there are three palaces $A_1,A_2,A_3\in \mac G$ of rank $j$, such that for $1\le i\le 3$,
		$A_i$ $\mac G$-adjoins some member of $D'$; and
		\item either $|D'|=1$, or $D'$ is a $\mac G$-village, or for each $j'\in \{k+1\LL \ell-1\}$ there are at most four palaces $A\in \mac G$ of rank $j'$ such that
		$A$ $\mac G$-adjoins a member of $D'$.
	\end{itemize}
	The first two bullets are clear.
	The third holds since no member of $D$ is $D$-peripheral except possibly the forts $X_{i_1}, X_{i_2}$, and no $D$-village is $D$-peripheral. Hence the only $D'$-peripheral $D'$-villages are those (at most three) that we added.
	
	For the fourth bullet, suppose that $X\in D'$ $\mac G$-adjoins some rebel $Y\notin D'$.
	If $X$ is a fort, then $X=X_i$ for some $i\in \{i_1\LL i_2\}$; and since $Y\notin D$, it follows that $i\in\{i_1,i_2\}$  and so $X$ is a leader fort.
	Now we assume that $X$ is a house;  let $B$ be the $\mac G$-village of $\mac G$ that contains $X$. Since $D'$
	is $\mac G$-village-closed, it follows that $B\subseteq D'$. Since $Y\notin D'$, it follows that $Y\notin B$, and so $B$ $\mac G$-adjoins $Y$ and hence $Y$
	is a fort. Choose $i\in \{i_1\LL i_2\}$ such that $B$ $\mac G$-adjoins $X_i$, with $i\notin \{i_1,i_2\}$ if possible.
	If $i\ne i_1,i_2$, then $X_i$ $\mac T$-semiadjoins $X_{i-1}, X_{i+1}$ and $Y$, contrary to \ref{contactgraph}.
	Hence we cannot choose $i\notin \{i_1,i_2\}$; so $B\not\subseteq D$, and therefore $B$ is a leader $\mac G$-village.
	This proves the fourth bullet.
	
	The fifth bullet holds since $D'$ is dangerous. Finally, for the sixth bullet, 
	suppose first that $D'\ne D$, and let $B$ be a $\mac G$-village in $D'\setminus D$. Let $j'\in \{k+1\LL \ell-1\}$.
	Since we added as few $\mac G$-villages to $D$ as possible to make a dangerous set, $D'\setminus B$ $\mac G$-adjoins at most two  members of $\mac G$ of rank $j'$; and since
	$B$ is not dangerous, $B$ also $\mac G$-adjoins at most two  members of $\mac G$ of rank $j'$. Hence $D'$ $\mac G$-adjoins at most four 
	members of $\mac G$ of rank $j'$ and the sixth bullet holds. So we may assume that $D'=D$, and so $D$ is dangerous.
	If $i_1\ne i_2$, then for every palace $A\in \mac G$, if $A$ $\mac G$-adjoins some member
	of $D$
	then it also $\mac G$-adjoins a member of one of $C(i_1+1,i_2),C(i_1,i_2-1)$, and since $C(i_1+1,i_2)$ and $C(i_1,i_2-1)$ are not dangerous (by the minimality of $i_2-i_1$), it follows that
	the sixth bullet holds. So we may assume that $i_1=i_2$. Since $D=D'$, it follows that $D=\{X_{i_1}\}$ and so $|D|=1$, and again the sixth bullet holds.
	
	Hence $D'$ is a cabal, a contradiction since $\mac G$ is a stable government.
	This proves (1).
	
	\bigskip
	
	Thus, $C$ is not dangerous, and so 
	for $k+1\le j\le \ell-1$, there are at most two palaces $A\in \mac G$ with rank $j$ such that $A$ $\mac G$-adjoins some member of $C$.
	Let $\mac Q$ be the set of all such palaces in $\mac G$; so $|\mac Q|\le 2(\ell-k-1)$. But since $C$ is a maximal
	$\mac G$-adjoin-connected set of rebels, it follows that for each $u\in \bd(\voro_{\mac G}(C))$ there exists $A\in \mac Q$ such that $u$ has a neighbour in $\voro_{\mac G}(A)$;
	and so by \ref{palacebound}, $\bd(\voro_{\mac G}(C))$ has quasi-size at most
	$$\left(8(\ell-k-1) 3^{\ell-k-1}\alpha, \beta+4d_0+1\right),$$
	and the theorem holds.
	This proves \ref{notmuchleft}.~\bbox

\end{document}
